\section{Language-Space Optimizers}
\label{sec:optimizers}

The improvement operator of \S\ref{sec:framework} can revise one output, write persistent memory, or update model parameters. Language-space optimization occupies the middle ground: an outer optimizer treats a prompt, context, playbook, or scaffold as an editable system parameter and improves it against a task metric. It is therefore not a fourth pillar. The optimized text persists across episodes like a learning artifact, but a frozen model consumes it through the context window. This semi-parametric view also clarifies the analogy to numerical optimization: critique supplies directional information, text is the parameter, an LLM edit is the update, and execution traces route credit through compound systems \citep{xiao2024verbalized,cheng2024trace}.

\subsection{Mechanism Families}
\label{subsec:optimizers-search}

Four families differ in the information supplied to the update. \emph{Score-driven search} proposes candidates and retains those with higher task scores; OPRO places previous candidates and scores in context so the optimizer can exploit its search history \citep{yang2024opro}. \emph{Textual-gradient methods} replace a scalar score with a diagnosis and propagate it to the responsible prompt or program component; TextGrad is the canonical realization, while trace-based methods extend the idea to compound graphs \citep{yuksekgonul2024textgrad}. \emph{Reflective and evolutionary methods} maintain diverse candidates, diagnose failures, and select complementary improvements; GEPA shows that richer diagnoses can achieve stronger sample efficiency than policy-gradient updates when rollouts are scarce \citep{agrawal2026gepa}. \emph{Hybrid methods} alternate text edits with weight updates, using prompt evolution to escape failures and then distilling the gain into the model; P$^2$O exemplifies this division of labor \citep{lu2026p2o}. Detailed systems and quantitative comparisons appear in the electronic supplement.

The families also assign credit at different resolutions. Score-only search credits a whole candidate, which is simple but expensive when pipelines contain many coupled prompts. Textual gradients use execution traces to localize blame to a component. Reflective populations trade precise attribution for diversity, retaining several explanations or solutions when no single update dominates every objective. Hybrids add a second choice: whether a discovered improvement should remain explicit in context or become procedural knowledge in weights. This credit unit, candidate, component, population, or parameter update, is often more consequential than the particular LLM used as optimizer.

\subsection{Selection Rules and Failure Modes}
\label{subsec:optimizers-tradeoffs}

Language-space optimization is preferred when weights or gradients are inaccessible, supervision is scarce enough that explanatory feedback must extract more information per rollout, or the learned artifact must remain inspectable and editable. Trace-aware textual gradients are especially useful when a compound pipeline needs component-level credit. Population methods help when several objectives or qualitatively different solutions must be retained. Weight updates remain preferable when abundant data amortizes training, the target knowledge is procedural, or deployment cannot afford a growing context. Hybrid methods fit the intermediate case, where an interpretable prompt discovers a behavior that must later be internalized.

The same properties create predictable failures. A reflective optimizer can misdiagnose the cause of an error and move in the wrong direction \citep{ma2024are}; wholesale context rewriting can collapse useful detail \citep{zhang2025agentic}; multi-objective critiques can dilute task focus \citep{darshan2026when}; and poisoned feedback directly poisons the update channel. Two contrast cases mark the definitional fence: text optimized by actual numeric gradients or by selection signals with no linguistic meaning is prompt optimization, but not verbal reinforcement learning \citep{das2024greater,wang2025evolving}. The verbal channel primarily buys sample efficiency, auditability, and an interface for structured credit, not a guarantee of higher attainable performance.
